\section{Energy Landscape Analysis}

\subsection{Group-level landscape characterization}

We characterized the group-level energy landscape derived from the population parameter vector $\boldsymbol{\eta}$ by enumerating all $2^d$ possible response patterns and computing their energies and equilibrium probabilities. To identify attractor states, we defined the neighborhood of each state as all states reachable by a single-bit flip (Hamming distance one). A state was classified as a local minimum if its energy was not greater than the energies of all its neighbors. We then partitioned the state space into basins of attraction using a deterministic steepest-descent algorithm: from each state, we iteratively moved to the neighbor with strictly lower energy until reaching a fixed point, which defined the basin attractor (Fig.~2A).

For each basin, we computed its probability mass as the sum of equilibrium probabilities over all states assigned to that basin (Fig.~2B). The basin probability mass reflects the expected occupancy under the equilibrium distribution and provides a natural ranking of basins by their statistical prominence. We summarized the landscape structure through scalar features including the number of local minima, the global minimum energy, the energy gap between the two lowest minima, the global basin mass, and the entropy of basin masses.

To characterize basin-to-basin dynamics, we constructed a coarse-grained transition matrix by aggregating single-bit-flip transitions under a Metropolis acceptance rule. Specifically, from each microstate we proposed a random single-bit flip uniformly among $d$ neighbors and accepted downhill moves with probability one and uphill moves with probability $\exp(-\Delta E)$, where $\Delta E$ is the energy difference. We then aggregated the resulting stationary probability fluxes across states to obtain a basin-level transition matrix $T_{\text{basin}}$ and basin stationary masses $\boldsymbol{\pi}_{\text{basin}}$ (Fig.~2C). The transition matrix was row-normalized to represent conditional transition probabilities between basins.

To visualize the energy barrier structure between major basins, we computed pairwise saddle energies using a minimax path criterion on the Hamming graph. For each pair of local minima, we identified the minimum-energy path connecting them and recorded the maximum energy along that path as the saddle energy. We then constructed a disconnectivity-like dendrogram by clustering basins based on their pairwise saddle energies (Fig.~2D). This representation reveals the hierarchical organization of energy barriers and helps identify metastable states separated by high barriers.

Finally, we visualized the full energy landscape by embedding the state space in two dimensions using a weighted graph layout where edges within the same basin were assigned higher weights than edges between basins. We overlaid the basin regions with color-coded basin assignments and energy contours, and annotated the positions of local minima and observed trajectory states (Fig.~2E).

\subsection{Individual-level landscape analysis}

For each subject, we constructed an individual energy landscape from their subject-specific parameter vector $\boldsymbol{\mu}_k$ using the same enumeration and basin assignment procedures. We visualized individual landscapes using the same weighted graph layout and basin coloring scheme as the group landscape (Fig.~3A, B).

To quantify the divergence between individual and group landscapes, we computed several distance measures on the full state-space probability distributions. The Jensen-Shannon divergence was computed as $JS(P_k, P_G) = \frac{1}{2} KL(P_k \| M) + \frac{1}{2} KL(P_G \| M)$, where $M = \frac{1}{2}(P_k + P_G)$ is the mixture distribution and $KL$ denotes the Kullback-Leibler divergence. We also computed asymmetric KL divergences $KL(P_k \| P_G)$ and $KL(P_G \| P_k)$, and the Pearson correlation coefficient between the energy vectors $E_k(s)$ and $E_G(s)$ across all states. Additionally, we computed differences in the number of local minima and in global minimum energies between individual and group landscapes (Fig.~3C).

\subsection{Trajectory analysis within landscapes}

To analyze how observed response trajectories relate to the energy landscapes, we mapped each observed response pattern at time $t$ to its corresponding state index in the enumerated state space. For each time point, we annotated the basin labels under both the group landscape ($b_G(t)$) and the individual landscape ($b_k(t)$), along with the corresponding energies and probabilities. We also computed pointwise differences $\Delta E(t) = E_G(s(t)) - E_k(s(t))$ and $\Delta \log P(t) = \log P_G(s(t)) - \log P_k(s(t))$ to quantify cross-landscape discrepancies at each time point.

We visualized trajectories by overlaying the observed state sequence on the energy landscape plots, with color-coding and line connections indicating the temporal progression through the state space (Fig.~3D, E for group landscape; Fig.~3F, G for individual landscape). To quantify the dynamics, we computed switching rates as the proportion of time steps where the basin label changed between consecutive observations. Specifically, the group switching rate was defined as the proportion of transitions where $b_G(t+1) \neq b_G(t)$, and the micro switching rate was defined as the proportion of transitions where $b_k(t+1) \neq b_k(t)$. These rates provide measures of basin stability and transition frequency within each landscape framework.
